
\section{Estimação de Estados}
\label{sec:estimacao}

Para realizar o seguimento autônomo de trajetórias, é necessário conhecer continuamente a posição, a orientação e a velocidade do veículo. Entretanto, as medições dos sensores embarcados estão sujeitas a ruídos, incertezas e diferentes frequências de atualização. Para combinar essas informações e obter estimativas consistentes dos estados do veículo, a plataforma VIDA utiliza um filtro de Kalman estendido (EKF), executado na Odroid, cujos resultados são disponibilizados ao sistema de navegação.

\subsection{Modelo de Estados}
\label{subsec:modelo_estados}

O estimador utiliza o modelo cinemático de bicicleta, apresentado na Seção~\ref{subsec:modelo_bicicleta}, para representar a evolução do movimento do veículo. Esse modelo estabelece as relações entre posição, orientação, velocidade longitudinal e ângulo de esterçamento, permitindo prever o comportamento do sistema ao longo do tempo.

O vetor de estados considerado pelo filtro é definido por:

\begin{equation}
\mathbf{x} =
\begin{bmatrix}
x_r & y_r & \theta & \delta & v & b_g
\end{bmatrix}^{T},
\label{eq:vetor_estados_ekf}
\end{equation}

em que $x_r$ e $y_r$ representam as coordenadas da posição do veículo em um referencial global, $\theta$ corresponde à sua orientação nesse referencial, $\delta$ ao ângulo de esterçamento, $v$ à velocidade longitudinal no eixo traseiro e $b_g$ ao \textit{bias} do giroscópio.

A inclusão do \textit{bias} do giroscópio no vetor de estados permite considerar erros sistemáticos associados à medição da velocidade angular. Dessa forma, além das variáveis que descrevem o movimento do veículo, o estimador contempla uma componente destinada à compensação de desvios presentes nas medições inerciais.

A partir desse modelo, a etapa de predição do EKF propaga os estados estimados e suas respectivas incertezas ao longo do tempo. As estimativas obtidas na etapa de predição são posteriormente atualizadas por meio das medições dos sensores, conforme os princípios apresentados na Subseção~\ref{subsec:fund_ekf}.

\subsection{Fusão Sensorial e Correção dos Estados}
\label{subsec:fusao_sensorial}

A estimação dos estados é realizada por meio de ciclos sucessivos de predição e correção, combinando o modelo cinemático de bicicleta com as medições dos sensores embarcados. Na implementação utilizada no VIDA, o EKF opera a uma frequência nominal de 50~Hz, propagando continuamente a posição e a orientação do veículo. As medições de velocidade longitudinal e ângulo de esterçamento são utilizadas na predição do movimento, enquanto o giroscópio e o receptor GNSS contribuem para a atualização da orientação e a correção dos desvios de posicionamento.

Na etapa de predição, a velocidade longitudinal é obtida pela média das medições dos \textit{encoders} dianteiros, enquanto o ângulo de esterçamento é fornecido pelo LVDT. Como a velocidade considerada corresponde àquela medida nas rodas dianteiras, o modelo incorpora uma compensação geométrica para determinar o deslocamento do centro do eixo traseiro, adotado como referência de posição do veículo. Dessa forma, a propagação das coordenadas é realizada por meio das equações:

\begin{equation}
\begin{aligned}
x_{k+1}^{-} &= x_k + v_k \cos(\theta_k)\Delta t, \\
y_{k+1}^{-} &= y_k + v_k \sin(\theta_k)\Delta t.
\end{aligned}
\label{eq:predicao_posicao_ekf}
\end{equation}

Para a atualização da orientação, a implementação utiliza prioritariamente a velocidade angular medida pelo giroscópio. Durante o movimento, essa medição é corrigida pelo \textit{bias} estimado, permitindo representar a rotação do veículo sem depender exclusivamente da relação cinemática entre velocidade e esterçamento. A evolução da orientação é calculada por:

\begin{equation}
\theta_{k+1}^{-} = \theta_k + (\omega_{m,k} - b_{g,k})\Delta t,
\label{eq:predicao_orientacao_ekf}
\end{equation}

em que $\omega_{m,k}$ corresponde à velocidade angular medida pelo giroscópio e $b_{g,k}$ representa o \textit{bias} estimado. Caso as medições inerciais deixem de estar disponíveis, o estimador utiliza o modelo cinemático de bicicleta (Equação~\ref{eq:modelo_cinematico}) como alternativa para calcular a velocidade angular, o que mantém a propagação da orientação durante interrupções temporárias na
comunicação com o sensor.

O \textit{bias} do giroscópio recebe tratamento específico para reduzir o acúmulo de erros na orientação estimada. Durante a inicialização, seu valor é determinado a partir da média das medições obtidas com o veículo parado. Posteriormente, quando os \textit{encoders} indicam velocidade inferior a $0{,}02~\mathrm{m/s}$ por um período mínimo de $0{,}5~\mathrm{s}$, o estimador considera que a velocidade angular real é nula. Nessa condição, a medição do giroscópio é utilizada na etapa de correção para atualizar o \textit{bias}, enquanto a integração da orientação é interrompida, evitando a acumulação de desvios durante as paradas.

Além da propagação dos estados, o filtro incorpora as medições de velocidade e esterçamento em sua etapa de correção, atualizando as respectivas incertezas. Embora essas grandezas sejam diretamente atualizadas pelas leituras dos sensores durante a predição, sua inclusão no processo de correção permite ajustar a matriz de covariância associada ao estado estimado.

A correção do posicionamento é realizada a partir das coordenadas fornecidas pelo receptor GNSS com correções RTK. A cada nova observação válida, disponibilizada nominalmente a 2~Hz, o filtro compara a posição medida com aquela prevista pelo modelo de movimento. As diferenças identificadas são utilizadas para atualizar a estimativa por meio do ganho de Kalman, considerando as incertezas associadas às medições GNSS. Dessa forma, os desvios acumulados durante a propagação são corrigidos periodicamente, enquanto as relações representadas na matriz de covariância permitem que essas atualizações também influenciem indiretamente outros estados, como a orientação.

A combinação dessas informações permite manter estimativas regulares de posição e orientação entre as atualizações do GNSS. Durante o deslocamento, a predição utiliza as medições de velocidade, esterçamento e velocidade angular para acompanhar o movimento do veículo, enquanto as observações de posicionamento limitam o acúmulo de desvios. Nas condições de parada, a atualização do \textit{bias} do giroscópio contribui para reduzir erros sistemáticos na estimação da orientação, proporcionando maior consistência às informações disponibilizadas pelo estimador.

O estimador utilizado neste trabalho baseia-se no filtro de Kalman estendido apresentado por \citet{DeLima2015}, posteriormente adaptado à plataforma VIDA e implementado como um nó ROS~2 para realizar a fusão das medições dos sensores embarcados e estimar os estados do veículo.
